
\documentclass[10pt,openany,a4paper]{article}
\usepackage{tikz, pgfplots, tkz-euclide,calc}
\usetikzlibrary{patterns,shapes.arrows,calc}
\usepackage{fancyhdr}
\usepackage{enumerate,enumitem}
\usepackage{cancel}
\usepackage{varwidth}

% TAMBAHKAN PACKAGE SENDIRI KALAU KURANG

\usepackage{geometry}
\geometry{
	left = 22mm,
	right = 22mm,
	top = 25mm,
	bottom = 25mm,
}
\usepackage{graphicx}
\usepackage{subcaption}
\usepackage{soul}
\usepackage{hyperref}
\usepackage{float}
\usepackage{setspace}
\usepackage{afterpage}
\usepackage{xcolor}
\usepackage[nosectionbib]{apacite}
\usepackage{lmodern}
\usepackage{mathtools}
\usepackage{upgreek}
\usepackage{siunitx}
% \usepackage{fontawesome5}
\usepackage{fontawesome6}
\usepackage{physics}
\usepackage{pdfpages}
\usepackage{hyperref}
\renewcommand{\baselinestretch}{1.3}
\usepackage{multicol}
\usepackage{hyperref,ragged2e,amsmath,multicol,gensymb,setspace,
fancyhdr,amsfonts,tikz,pgfplots,nccmath,enumerate,verbatim,amsthm,physics}
\usepgfplotslibrary{polar,fillbetween}
% Disable tikz externalization to avoid requiring -shell-escape
\usepgfplotslibrary{external}
\tikzexternalize[prefix=tikz/]
\tikzset{reuse path/.code={\pgfsyssoftpath@setcurrentpath{#1}}}
\usetikzlibrary{shapes.geometric,arrows}
\newcommand\mylog[1]{\mathop{{}^{#1}\mathrm{log}}}
\pgfplotsset{compat=1.15}
\pgfplotsset{my style/.append style={axis x line=middle, axis y line=
middle, xlabel={$x$}, ylabel={$y$}, axis equal }}
\usepackage{etoolbox}
% remove stray \makeatother
\newtheorem{defn}{Definisi}[section]
\newtheorem{teo}[defn]{Teorema}
\newtheorem{thm}{Teorema}[section]
\newtheorem{lemma}{Lemma}
\newtheorem{lemmas}[defn]{Lemma}
\newtheorem{cor}[defn]{Akibat}
\newtheorem{asumsi}[defn]{Asumsi}
\theoremstyle{definition}
\newtheorem{con}{Contoh}[section]
\theoremstyle{definition}
\theoremstyle{plain}
\newtheorem{prop}{Proposisi}[section]
\renewcommand{\proofname}{Bukti}
\renewcommand{\thethm}{\arabic{chapter}.\arabic{thm}}

\newcommand{\normlp}[1]{\left\|#1\right\|_{L^p_{\lambda}(0,T;H)}}% Fungsi norm (||x||)
\newcommand{\R}{\mathbb{R}}
\renewcommand{\natural}{\mathbb{N}}
\newcommand{\Xpl}{X^R_{p,\lambda}}
\newcommand{\Xpls}{X^R_{p,\lambda^*}}
\newcommand{\qpl}{q^R_{p,\lambda}}
\newcommand{\qpls}{q^R_{p,\lambda^*}}
\pagestyle{fancy}
\fancyhf{}
\lhead{Halaman \thepage}
% \cfoot{\faXTwitter\ \href{https://x.com/ahmadzakiyudin_/}{@ahmadzakiyudin\_} | \faInstagram\ \href{https://instagram.com/ahmadzakiyudin_/}{@ahmadzakiyudin\_} | \faTiktok\ \href{https://tiktok.com/@ahmathz/}{@ahmathz}}
% \rhead{Pembahasan Soal ETS \\ Kalkulus 1 ITS 2025/2026}
\hypersetup{
    colorlinks=true,
    linkcolor=blue,
    filecolor=blue,      
    urlcolor=blue,
}
\newenvironment{jawab}[2][tes]{
% \fbox{\begin{minipage}{14.65cm}
\textbf{Kata kunci materi:} #1.\\
\textbf{Penyelesaian:} #2
% \end{minipage}}
}{}

\newenvironment{jawab2}[2][tes]{
\fbox{\begin{minipage}{14.65cm}
 #2
\end{minipage}}}{}

\newenvironment{jawaban}[2][tes]{
\colorlet{oldcolor}{.}
        \color{red}
        \setlength{\fboxrule}{1pt}\boxed{\color{oldcolor}#2}\color{oldcolor}.
}{}

\newenvironment{jawaban2}[2][tes]{
\colorlet{oldcolor}{.}
        \color{red}
        \setlength{\fboxrule}{1pt}\fbox{\color{oldcolor}#2}\color{oldcolor}.
}{}

\newenvironment{penting}[2][tes]{
\colorlet{oldcolor}{.}
\color{blue}
\fbox{\color{oldcolor}#2}\color{oldcolor}
}{}

\begin{document}
% \begin{titlepage}
%     \vspace*{\fill}
%     \begin{center}
%       \Huge {PEMBAHASAN SOAL EAS \\ KALKULUS II \\ TAHUN 2024/2025}\\[0.4 cm]
%       \huge {Ahmad Hisbu Zakiyudin}\\[0.2 cm]
%       \Large {(\href{https://instagram.com/ahmadzakiyudin_/}{@ahmadzakiyudin\_})}
%     \end{center}
%     \vspace*{\fill}
%   \end{titlepage}
% !TEX root = ../../BookTemplate.tex
%%%%%%%%%%%%%%%%%%%%%%%%%%%%%%%%%%%%%%%%%%%%%%%%%%%%%%%%%%%%%%%%%%%%%%%%%%%%%%%%%%
% \begin{titlepage}
%     \raggedleft	
% 	\hspace{.025\textwidth}
% 	\parbox[b]{\textwidth}{
%     \vspace{1.5cm}
% 		{\Huge\bfseries Pembahasan Soal EAS}                \\[20pt]
% 		{\huge\textit\ Kalkulus II}\\ \\[10pt]
%         {\Large\textsc{Tahun 2024/2025}}}
% \vfill
% \rule{1pt}{.14\textheight}
% \hspace{.025\textwidth}
%     \raisebox{10pt}{\parbox[b]{.95\textwidth}{ 
%         {\Large\textsc{Ahmad Hisbu Zakiyudin}}              \\[10pt]
%         \large{\href{https://instagram.com/ahmadzakiyudin_/}{@ahmadzakiyudin\_}}               \\[25pt]
%         {\large 5 Juni 2026}}}
% \end{titlepage}
\begin{titlepage}
	\centering % Center everything on the title page
	% \scshape % Use small caps for all text on the title page
	\vspace*{1.5\baselineskip} % White space at the top of the page
% ===================
%	Title Section 	
% ===================

	\rule{13cm}{1.6pt}\vspace*{-\baselineskip}\vspace*{2pt} % Thick horizontal rule
	\rule{13cm}{0.4pt} % Thin horizontal rule
	
		\vspace{0.75\baselineskip} % Whitespace above the title
% ========== Title ===============	
	{	\Huge Soal Tryout ETS\\ 
			\vspace{4mm}
		Kalkulus 1 \\ \vspace{2mm} \huge{Tahun 2026/2027}\\ \vspace{2mm} \Large{(Paket 2)}\\	}
% ======================================
		\vspace{0.75\baselineskip} % Whitespace below the title
	\rule{13cm}{0.4pt}\vspace*{-\baselineskip}\vspace{3.2pt} % Thin horizontal rule
	\rule{13cm}{1.6pt} % Thick horizontal rule
	
		\vspace{1.75\baselineskip} % Whitespace after the title block
% =================
%	Information	
% =================
{\large Oleh: Ahmad Hisbu Zakiyudin}
	% {\large Oleh: Ahmad Hisbu Zakiyudin \\
	% 	\vspace*{1.2\baselineskip}
	% \href{https://instagram.com/ahmadzakiyudin_/}{@ahmadzakiyudin\_}} \\
	\vfill
	\vspace{1mm}
\huge{Selamat Belajar!}
\end{titlepage}
\makeatletter
\setlength\fboxsep{3pt}
% ensure enough header height for fancyhdr
\setlength{\headheight}{24pt}
% Kop tiap kelompok: #1 kelompok, #2 kelas, #3 departemen, #4 jadwal kuliah, #5 sesi
\newcommand{\kopsoal}[5]{%
 \begin{center}
	{\underline{\textbf{\MakeUppercase{Tryout ETS Kalkulus 1 Tahun 2026/2027}}}}
    \end{center}

    \begin{center}
	\begin{tabular}{lcl}
		Mata kuliah/SKS & : & Kalkulus 1 ( SM234101 ) / 3 SKS\\
		Kelompok & : & #1\\
		Kelas & : & #2\\
		Departemen & : & #3\\
		Jadwal kuliah & : & #4\\
		Sesi & : & #5\\
		Waktu & : & 100 menit\\
		Sifat & : & Tertutup\\
  ~
	\end{tabular}

    \noindent\rule{\textwidth}{2.pt}
        \end{center}}

% Kop kelompok IUP (bahasa Inggris): #1 kelompok, #2 kelas, #3 departemen, #4 jadwal kuliah, #5 sesi
\newcommand{\kopsoaliup}[5]{%
 \begin{center}
	{\underline{\textbf{\MakeUppercase{Tryout ETS Kalkulus 1 Tahun 2026/2027}}}}
    \end{center}

    \begin{center}
	\begin{tabular}{lcl}
		Course / Credit & : & Kalkulus 1 ( SM234101 ) / 3 SKS\\
		Group & : & #1\\
		Class & : & #2\\
		Department & : & #3\\
		Lecture schedule & : & #4\\
		Session & : & #5\\
		Time & : & 100 minutes\\
		Type & : & Closed\\
  ~
	\end{tabular}

    \noindent\rule{\textwidth}{2.pt}
        \end{center}}

% Kop tryout Matematika 2 SKS: #1 kelompok, #2 kelas, #3 departemen, #4 jadwal kuliah, #5 sesi
\newcommand{\kopsoalmat}[5]{%
 \begin{center}
	{\underline{\textbf{\MakeUppercase{Tryout ETS Matematika Tahun 2026/2027}}}}
    \end{center}

    \begin{center}
	\begin{tabular}{lcl}
		Mata kuliah/SKS & : & Matematika ( SM234152 ) / 2 SKS\\
		Kelompok & : & #1\\
		Kelas & : & #2\\
		Departemen & : & #3\\
		Jadwal kuliah & : & #4\\
		Sesi & : & #5\\
		Waktu & : & 100 menit\\
		Sifat & : & Tertutup\\
  ~
	\end{tabular}

    \noindent\rule{\textwidth}{2.pt}
        \end{center}}

% Kop tryout Matematika 2 SKS IUP (bahasa Inggris)
\newcommand{\kopsoalmatiup}[5]{%
 \begin{center}
	{\underline{\textbf{\MakeUppercase{Tryout ETS Matematika Tahun 2026/2027}}}}
    \end{center}

    \begin{center}
	\begin{tabular}{lcl}
		Course / Credit & : & Matematika ( SM234152 ) / 2 SKS\\
		Group & : & #1\\
		Class & : & #2\\
		Department & : & #3\\
		Lecture schedule & : & #4\\
		Session & : & #5\\
		Time & : & 100 minutes\\
		Type & : & Closed\\
  ~
	\end{tabular}

    \noindent\rule{\textwidth}{2.pt}
        \end{center}}

%%%%%%%%%%%%%%%%%%%%%%%%%%%%%%%%%%%%%%%%%%%%%%%%%%%%%%%%%%%%%%%%%%%%%%%%%%%%%%
% Kal-5 | Paket 2: subbab per nomor 1.1, 2.3, 2.4, 3.3, 4.2
% Kisi-kisi hari Senin: subbab 1.1, 1.5, 2.3, 2.4, 3.1, 3.3, 4.2, 4.4
\kopsoal{Kal-5}{22-27}{T. Fisika, T. Material dan Metalurgi, Rkp}{Senin, Rabu; 07.00-09.30}{1}
\begin{enumerate}
    \item \textbf{[Subbab 1.1]} Dapatkan himpunan penyelesaian pertidaksamaan
    \begin{align*}
        \frac{2x-1}{x+3}\leq 1.
    \end{align*}

    \item \textbf{[Subbab 2.3]} Gambarkan kurva dasar $y=\sqrt{x}$. Berdasarkan kurva dasar tersebut, gunakan translasi (pergeseran) untuk menggambar kurva $y=\sqrt{x-2}+1$ (Jelaskan ilustrasinya).

    \item \textbf{[Subbab 2.4]} Diberikan fungsi $f(x)=\dfrac{3x-2}{x+1}$.
    \begin{enumerate}
        \item Tentukan domain dan range dari $f$.
        \item Dapatkan $f^{-1}(x)$ beserta domainnya.
    \end{enumerate}

    \item \textbf{[Subbab 3.3]} Hitunglah $\displaystyle\lim_{x\to-\infty}\frac{2x+1}{\sqrt{x^2+3}}$.

    \item \textbf{[Subbab 4.2]} Dengan menggunakan definisi turunan, dapatkan $f'(x)$ dari $f(x)=x^2+3x$. Kemudian dapatkan persamaan garis singgung kurva $y=f(x)$ di $x=1$.
\end{enumerate}

\newpage

%%%%%%%%%%%%%%%%%%%%%%%%%%%%%%%%%%%%%%%%%%%%%%%%%%%%%%%%%%%%%%%%%%%%%%%%%%%%%%
% Kal-6 | Paket 2: subbab per nomor 1.1, 2.3, 2.4, 3.3, 4.2
% Kisi-kisi hari Senin: subbab 1.1, 1.5, 2.3, 2.4, 3.1, 3.3, 4.2, 4.4
\kopsoal{Kal-6}{28-31}{T. Perkapalan, T. Transportasi Laut, T. Bangunan Lepas Pantai}{Senin, Rabu; 07.00-09.30}{1}
\begin{enumerate}
    \item \textbf{[Subbab 1.1]} Selesaikan pertidaksamaan berikut
    \begin{align*}
        \frac{x^2-3x-4}{x^2-9}\leq 0.
    \end{align*}

    \item \textbf{[Subbab 2.3]} Gambarkan sketsa grafik $f(x)=|x-3|$ dan $g(x)=|x+1|$ pada satu bidang koordinat. Tentukan titik potong kedua grafik tersebut.

    \item \textbf{[Subbab 2.4]} Diberikan fungsi $f(x)=2+\sqrt{5-x}$ dengan $x\leq 5$.
    \begin{enumerate}
        \item Tentukan range dari $f$.
        \item Dapatkan $f^{-1}(x)$ beserta domainnya.
        \item Sketsa grafik $f(x)$ dan $f^{-1}(x)$ pada satu bidang koordinat.
    \end{enumerate}

    \item \textbf{[Subbab 3.3]} Hitunglah $\displaystyle\lim_{x\to+\infty}\frac{3x^2-x+2}{1-2x^2}$.

    \item \textbf{[Subbab 4.2]} Tunjukkan bahwa $f(x)=\begin{cases}
        x^2+3, & x\leq 1\\
        2x+2, & x>1
    \end{cases}$ kontinu dan dapat diturunkan di $x=1$.
\end{enumerate}

\newpage

%%%%%%%%%%%%%%%%%%%%%%%%%%%%%%%%%%%%%%%%%%%%%%%%%%%%%%%%%%%%%%%%%%%%%%%%%%%%%%
% Kal-iup-1 | Paket 2: subbab per nomor 1.1, 2.3, 2.4, 3.1, 4.2
% Kisi-kisi hari Senin: subbab 1.1, 1.5, 2.3, 2.4, 3.1, 3.3, 4.2, 4.4
\kopsoaliup{Kal-iup-1}{61, 68}{Chemical, Environmental, Electrical Engineering}{Monday, Wednesday; 07.00-09.30}{1}
\begin{enumerate}
    \item \textbf{[Subbab 1.1]} Solve the following inequality.
    \begin{align*}
        \frac{x+5}{x-1}\geq 2.
    \end{align*}

    \item \textbf{[Subbab 2.3]} Sketch the base curve $y=\dfrac{1}{x}$. Based on that base curve, use translation to sketch the curve $y=\dfrac{1}{x-2}+1$. (Explain the process).

    \item \textbf{[Subbab 2.4]} Given a function $f(x)=x^2-6x+5$ with $x\geq 3$. Determine:
    \begin{enumerate}
        \item[a)] The range of $f$;
        \item[b)] $f^{-1}(x)$ and its domain.
    \end{enumerate}

    \item \textbf{[Subbab 3.1]} Given $f(x)=\begin{cases}
        x+1, & x<1\\
        5-x^2, & x\geq 1
    \end{cases}$.
    \begin{enumerate}
        \item[a.] Sketch the graph of $f$.
        \item[b.] Find $\displaystyle\lim_{x\to 1^-}f(x)$ and $\displaystyle\lim_{x\to 1^+}f(x)$. Does $\displaystyle\lim_{x\to 1}f(x)$ exist? Explain.
    \end{enumerate}

    \item \textbf{[Subbab 4.2]} Using the definition of derivative, find $f'(x)$ of $f(x)=\sqrt{x+3}$.
\end{enumerate}

\newpage

%%%%%%%%%%%%%%%%%%%%%%%%%%%%%%%%%%%%%%%%%%%%%%%%%%%%%%%%%%%%%%%%%%%%%%%%%%%%%%
% Kal-iup-4 | Paket 2: subbab per nomor 1.1, 2.3, 2.4, 3.1, 4.2
% Kisi-kisi hari Senin: subbab 1.1, 1.5, 2.3, 2.4, 3.1, 3.3, 4.2, 4.4
\kopsoaliup{Kal-iup-4}{65, 66, 67}{Statistics, Marine Engineering (JD), Informatics}{Monday, Wednesday; 07.00-09.30}{1}
\begin{enumerate}
    \item \textbf{[Subbab 1.1]} Solve the following inequality.
    \begin{align*}
        3\leq x^2+2x<8.
    \end{align*}

    \item \textbf{[Subbab 2.3]} Sketch the base curve $y=|x|$. Based on that base curve, use translation to sketch the curve $y=|x+2|-3$. (Explain the process).

    \item \textbf{[Subbab 2.4]} Given a function $f(x)=\dfrac{x-1}{2x+3}$. Determine:
    \begin{enumerate}
        \item[a)] The domain and range of $f$;
        \item[b)] $f^{-1}(x)$ and its domain.
    \end{enumerate}

    \item \textbf{[Subbab 3.1]} Given $f(x)=\begin{cases}
        2-x, & x<0\\
        x^2, & x\geq 0
    \end{cases}$.
    \begin{enumerate}
        \item[a.] Sketch the graph of $f$.
        \item[b.] Find $\displaystyle\lim_{x\to 0^-}f(x)$, $\displaystyle\lim_{x\to 0^+}f(x)$, and $\displaystyle\lim_{x\to 1}f(x)$ (if they exist).
    \end{enumerate}

    \item \textbf{[Subbab 4.2]} Using the definition of derivative, find $f'(x)$ of $f(x)=\dfrac{3}{x}$. Then, find the equation of the tangent line to the curve $y=f(x)$ at $x=1$.
\end{enumerate}

\newpage

%%%%%%%%%%%%%%%%%%%%%%%%%%%%%%%%%%%%%%%%%%%%%%%%%%%%%%%%%%%%%%%%%%%%%%%%%%%%%%
% Kal-9 | Paket 2: subbab per nomor 1.1, 2.3, 2.4, 3.1, 4.2
% Kisi-kisi hari Senin: subbab 1.1, 1.5, 2.3, 2.4, 3.1, 3.3, 4.2, 4.4
\kopsoal{Kal-9}{44-48}{T. Sistem Perkapalan, T. Kelautan, T. Komputer}{Senin, Rabu; 12.30-15.00}{3}
\begin{enumerate}
    \item \textbf{[Subbab 1.1]} Dapatkan himpunan penyelesaian dari
    \begin{align*}
        \frac{x+1}{x-2}\geq\frac{x-3}{x+1}.
    \end{align*}

    \item \textbf{[Subbab 2.3]} Diberikan suatu fungsi $f(x)=1+\sqrt{x+2}$ dan garis $\ell$ yang melalui titik $(-2,1)$ dan $(7,4)$.
    \begin{enumerate}
        \item Buatlah sketsa kurva $f(x)$ dan garis $\ell$ pada satu bidang koordinat.
        \item Tentukan titik potong antara kurva $f(x)$ dan garis $\ell$.
    \end{enumerate}

    \item \textbf{[Subbab 2.4]} Diberikan $f(x)=2-\sqrt[3]{x+1}$.
    \begin{enumerate}
        \item Dapatkan $f^{-1}(x)$ beserta domainnya.
        \item Sketsa $f(x)$ dan $f^{-1}(x)$ pada satu bidang koordinat.
    \end{enumerate}

    \item \textbf{[Subbab 3.1]} Diberikan $g(x)=\dfrac{x^2+x}{|x|}$. Dapatkan $\displaystyle\lim_{x\to 0^-}g(x)$ dan $\displaystyle\lim_{x\to 0^+}g(x)$. Apakah $\displaystyle\lim_{x\to 0}g(x)$ ada? Jelaskan.

    \item \textbf{[Subbab 4.2]} Dengan menggunakan definisi turunan, dapatkan $f'(x)$ dari $f(x)=\dfrac{1}{x-2}$. Kemudian tentukan kemiringan garis singgung kurva $y=f(x)$ di $x=3$.
\end{enumerate}

\newpage

%%%%%%%%%%%%%%%%%%%%%%%%%%%%%%%%%%%%%%%%%%%%%%%%%%%%%%%%%%%%%%%%%%%%%%%%%%%%%%
% Kal-11 | Paket 2: subbab per nomor 1.1, 2.3, 2.4, 3.3, 4.2
% Kisi-kisi hari Senin: subbab 1.1, 1.5, 2.3, 2.4, 3.1, 3.3, 4.2, 4.4
\kopsoal{Kal-11}{51-54}{Biologi, Bioteknologi, Teknologi Informasi}{Senin, Rabu; 12.30-15.00}{3}
\begin{enumerate}
    \item \textbf{[Subbab 1.1]} Tentukan semua nilai $x$ yang memenuhi pertidaksamaan
    \begin{align*}
        \frac{x^3-4x}{x^2-1}\geq 0.
    \end{align*}

    \item \textbf{[Subbab 2.3]} Gambarkan sketsa grafik fungsi berikut, kemudian tentukan range-nya.
    \begin{align*}
        f(x)=\begin{cases}
            4-x^2, & x\leq 1\\
            2x+1, & x>1
        \end{cases}
    \end{align*}

    \item \textbf{[Subbab 2.4]} Diberikan fungsi $f(x)=x^2+4x$ dengan $x\leq -2$.
    \begin{enumerate}
        \item Dapatkan range dari $f(x)$.
        \item Dapatkan $f^{-1}(x)$ beserta domainnya.
        \item Sketsa grafik $f(x)$ dan $f^{-1}(x)$ pada satu bidang koordinat.
    \end{enumerate}

    \item \textbf{[Subbab 3.3]} Hitunglah $\displaystyle\lim_{x\to-\infty}\left(\sqrt{x^2+2x}+x\right)$.

    \item \textbf{[Subbab 4.2]} Diberikan $f(x)=\begin{cases}
        x^2, & x\leq 2\\
        4x-4, & x>2
    \end{cases}$.\\
    Selidiki apakah $f$ dapat diturunkan di $x=2$. Jika dapat, tentukan $f'(2)$.
\end{enumerate}

\newpage

%%%%%%%%%%%%%%%%%%%%%%%%%%%%%%%%%%%%%%%%%%%%%%%%%%%%%%%%%%%%%%%%%%%%%%%%%%%%%%
% Kal-12 | Paket 2: subbab per nomor 1.1, 2.3, 2.4, 3.1, 4.2
% Kisi-kisi hari Senin: subbab 1.1, 1.5, 2.3, 2.4, 3.1, 3.3, 4.2, 4.4
\kopsoal{Kal-12}{55-60}{Aktuaria, Teknik Informatika, RPL, RKA}{Senin, Rabu; 12.30-15.00}{3}
\begin{enumerate}
    \item \textbf{[Subbab 1.1]} Tentukan penyelesaian dari pertidaksamaan
    \begin{align*}
        \frac{x^2-x-6}{x^2+x+1}<0.
    \end{align*}

    \item \textbf{[Subbab 2.3]} Gambarkan kurva dasar $y=x^3$. Berdasarkan kurva dasar tersebut, gunakan translasi (pergeseran) untuk menggambar kurva $y=(x+1)^3-2$ (Jelaskan ilustrasinya).

    \item \textbf{[Subbab 2.4]} Diberikan suatu fungsi $f(x)=1-\sqrt{9-x^2}$ dengan domain $[0,3]$.
    \begin{enumerate}
        \item Tentukan range dari $f$.
        \item Dapatkan $f^{-1}(x)$ beserta domainnya.
        \item Sketsa grafik $f(x)$ dan $f^{-1}(x)$ pada satu bidang koordinat.
    \end{enumerate}

    \item \textbf{[Subbab 3.1]} Diberikan $f(x)=\begin{cases}
        ax+3, & x<-1\\
        x^2+a, & x\geq -1
    \end{cases}$.\\
    Tentukan nilai $a$ agar $\displaystyle\lim_{x\to -1}f(x)$ ada.

    \item \textbf{[Subbab 4.2]} Dengan menggunakan definisi turunan, dapatkan $f'(x)$ dari $f(x)=x^3$. Kemudian dapatkan persamaan garis singgung kurva $y=f(x)$ di $x=1$.
\end{enumerate}

\newpage

%%%%%%%%%%%%%%%%%%%%%%%%%%%%%%%%%%%%%%%%%%%%%%%%%%%%%%%%%%%%%%%%%%%%%%%%%%%%%%
% Kal-7 | Paket 2: subbab per nomor 1.1, 2.3, 2.4, 3.1, 4.2
% Kisi-kisi hari Senin: subbab 1.1, 1.5, 2.3, 2.4, 3.1, 3.3, 4.2, 4.4
\kopsoal{Kal-7}{32-37}{T. Sipil, T. Geomatika, T. Geofisika, T. Pertambangan}{Senin, Rabu; 15.30-17.10}{4}
\begin{enumerate}
    \item \textbf{[Subbab 1.1]} Dapatkan penyelesaian dari pertidaksamaan
    \begin{align*}
        \frac{2x+3}{x-1}\geq 3.
    \end{align*}

    \item \textbf{[Subbab 2.3]} Gambarkan sketsa grafik fungsi
    \begin{align*}
        g(x)=\begin{cases}
            2-x, & x<1\\
            1+\sqrt{x-1}, & x\geq 1
        \end{cases}.
    \end{align*}

    \item \textbf{[Subbab 2.4]} Diberikan fungsi $f(x)=\dfrac{4}{x-2}+1$.
    \begin{enumerate}
        \item Tentukan domain dan range dari $f$.
        \item Dapatkan $f^{-1}(x)$ beserta domainnya.
    \end{enumerate}

    \item \textbf{[Subbab 3.1]} Diberikan $h(x)=\begin{cases}
        x+2, & x\leq 1\\
        x^2+1, & x>1
    \end{cases}$.
    \begin{enumerate}
        \item Sketsa grafik $h$.
        \item Dapatkan $\displaystyle\lim_{x\to 1^-}h(x)$ dan $\displaystyle\lim_{x\to 1^+}h(x)$. Apakah $\displaystyle\lim_{x\to 1}h(x)$ ada?
    \end{enumerate}

    \item \textbf{[Subbab 4.2]} Dengan menggunakan definisi turunan, dapatkan $f'(x)$ dari $f(x)=2x^2-x+1$. Kemudian dapatkan persamaan garis singgung kurva $y=f(x)$ di $x=1$.
\end{enumerate}

\newpage

%%%%%%%%%%%%%%%%%%%%%%%%%%%%%%%%%%%%%%%%%%%%%%%%%%%%%%%%%%%%%%%%%%%%%%%%%%%%%%
% Kal-8 | Paket 2: subbab per nomor 1.1, 1.5, 2.4, 3.1, 4.2
% Kisi-kisi hari Senin: subbab 1.1, 1.5, 2.3, 2.4, 3.1, 3.3, 4.2, 4.4
\kopsoal{Kal-8}{38-43}{Matematika, Statistika, Sains Data}{Senin, Rabu; 15.30-17.10}{4}
\begin{enumerate}
    \item \textbf{[Subbab 1.1]} Dapatkan himpunan penyelesaian pertidaksamaan
    \begin{align*}
        \frac{x+2}{x-3}\leq\frac{x}{x+1}.
    \end{align*}

    \item \textbf{[Subbab 1.5]} Diberikan persamaan setengah lingkaran
    \begin{align*}
        x=2-\sqrt{8+2y-y^2}.
    \end{align*}
    \begin{enumerate}
        \item Dapatkan pusat dan jari-jari lingkaran.
        \item Sketsa grafiknya.
    \end{enumerate}

    \item \textbf{[Subbab 2.4]} Diberikan $f(x)=3+2\sqrt[3]{x-1}$.
    \begin{enumerate}
        \item Dapatkan $f^{-1}(x)$.
        \item Sketsa grafik $f(x)$ dan $f^{-1}(x)$ pada satu bidang koordinat.
    \end{enumerate}

    \item \textbf{[Subbab 3.1]} Diberikan $f(x)=\begin{cases}
        k^2x-1, & x<2\\
        k(x+1)+1, & x\geq 2
    \end{cases}$.\\
    Tentukan semua nilai $k$ yang memenuhi agar $\displaystyle\lim_{x\to 2}f(x)$ ada.

    \item \textbf{[Subbab 4.2]} Tunjukkan bahwa $f(x)=\begin{cases}
        x^2-1, & x<1\\
        1-x, & x\geq 1
    \end{cases}$ kontinu tetapi tidak dapat diturunkan di $x=1$.
\end{enumerate}

\newpage

%%%%%%%%%%%%%%%%%%%%%%%%%%%%%%%%%%%%%%%%%%%%%%%%%%%%%%%%%%%%%%%%%%%%%%%%%%%%%%
% Kal-iup-2 | Paket 2: subbab per nomor 1.1, 2.3, 2.4, 3.1, 4.2
% Kisi-kisi hari Senin: subbab 1.1, 1.5, 2.3, 2.4, 3.1, 3.3, 4.2, 4.4
\kopsoaliup{Kal-iup-2}{62, 69}{Mechanical, Industrial, Materials \& Metallurgical Engineering}{Monday, Wednesday; 15.30-17.10}{4}
\begin{enumerate}
    \item \textbf{[Subbab 1.1]} Find the solution of the following inequality
    \begin{align*}
        \frac{x^2-4}{x^2+3x}\geq 0.
    \end{align*}

    \item \textbf{[Subbab 2.3]} Sketch the base curve $y=\sqrt{x}$. Based on that base curve, use translation to sketch the curve $y=\sqrt{x+4}-1$. (Explain the process).

    \item \textbf{[Subbab 2.4]} Given a function $f(x)=\sqrt{2x+6}-1$. Determine:
    \begin{enumerate}
        \item[a)] The domain and range of $f$;
        \item[b)] $f^{-1}(x)$ and its domain.
    \end{enumerate}

    \item \textbf{[Subbab 3.1]} Given $f(x)=\begin{cases}
        3-x, & x<2\\
        x^2-2, & x\geq 2
    \end{cases}$.
    \begin{enumerate}
        \item[a.] Sketch the graph of $f$.
        \item[b.] Find $\displaystyle\lim_{x\to 2^-}f(x)$ and $\displaystyle\lim_{x\to 2^+}f(x)$. Does $\displaystyle\lim_{x\to 2}f(x)$ exist?
    \end{enumerate}

    \item \textbf{[Subbab 4.2]} Show that $f(x)=\begin{cases}
        x^2+2, & x\leq 1\\
        3x, & x>1
    \end{cases}$ is continuous but not differentiable at $x=1$.
\end{enumerate}

\newpage

%%%%%%%%%%%%%%%%%%%%%%%%%%%%%%%%%%%%%%%%%%%%%%%%%%%%%%%%%%%%%%%%%%%%%%%%%%%%%%
% Kal-3 | Paket 2: subbab per nomor 1.2, 1.4, 2.1, 3.4, 4.2
% Kisi-kisi hari Selasa: subbab 1.2, 1.4, 2.1, 2.2, 3.2, 3.4, 4.2, 4.3
\kopsoal{Kal-3}{13-16}{Kimia, Saik, Telkom}{Selasa, Kamis; 07.00-09.30}{6}
\begin{enumerate}
    \item \textbf{[Subbab 1.2]} Tentukan himpunan penyelesaian dari persamaan
    \begin{align*}
        |x-2|+|x+1|=5.
    \end{align*}

    \item \textbf{[Subbab 1.4]} Garis $\ell$ melalui titik $(2,3)$ dan mempunyai gradien $-\frac{1}{2}$.
    \begin{enumerate}
        \item Dapatkan persamaan garis $\ell$.
        \item Jika garis $\ell$ memotong sumbu-$x$ di titik $A$ dan sumbu-$y$ di titik $B$, tentukan luas segitiga $OAB$ dengan $O$ titik asal.
    \end{enumerate}

    \item \textbf{[Subbab 2.1]} Diberikan suatu fungsi $f(x)=\left|x^2-4x+3\right|$.
    \begin{enumerate}
        \item Nyatakan $f(x)$ dalam bentuk fungsi sepotong-sepotong tanpa menggunakan tanda mutlak.
        \item Tentukan domain dan range dari $f(x)$.
    \end{enumerate}

    \item \textbf{[Subbab 3.4]} Dapatkan nilai $x$ di mana $f(x)=\dfrac{x^2-x-6}{x-3}$ diskontinu dan jelaskan apakah diskontinuitas tersebut dapat dihilangkan. Jika dapat, definisikan nilai fungsi di titik tersebut agar $f$ kontinu.

    \item \textbf{[Subbab 4.2]} Dengan menggunakan definisi turunan, dapatkan $f'(x)$ dari $f(x)=x^2-x$. Kemudian dapatkan persamaan garis singgung kurva $y=f(x)$ di $x=2$.
\end{enumerate}

\newpage

%%%%%%%%%%%%%%%%%%%%%%%%%%%%%%%%%%%%%%%%%%%%%%%%%%%%%%%%%%%%%%%%%%%%%%%%%%%%%%
% Kal-4 | Paket 2: subbab per nomor 1.2, 1.4, 2.2, 3.2, 4.2
% Kisi-kisi hari Selasa: subbab 1.2, 1.4, 2.1, 2.2, 3.2, 3.4, 4.2, 4.3
\kopsoal{Kal-4}{17-21}{T. Mesin, Teknik Sistem dan Industri}{Selasa, Kamis; 07.00-09.30}{6}
\begin{enumerate}
    \item \textbf{[Subbab 1.2]} Tentukan semua nilai $x$ yang memenuhi
    \begin{align*}
        1\leq|2x-3|<5.
    \end{align*}

    \item \textbf{[Subbab 1.4]} Diketahui garis $\ell_1$ melalui titik $(2,1)$ dan $(-2,-1)$. Jika $\ell_2$ adalah garis yang melalui titik $(1,-2)$ dan sejajar dengan garis $\ell_1$, dapatkan persamaan garis $\ell_2$ dan jarak titik $(2,1)$ ke garis $\ell_2$.

    \item \textbf{[Subbab 2.2]} Diberikan $f(x)=\sqrt{x+5}$ dan $g(x)=x^2-9$.
    \begin{enumerate}
        \item Dapatkan domain $f$ dan $g$.
        \item Dapatkan $(f\circ g)(x)$ beserta domainnya.
    \end{enumerate}

    \item \textbf{[Subbab 3.2]} Hitunglah limit berikut.
    \begin{enumerate}
        \item $\displaystyle\lim_{x\to 0}\frac{\sqrt{x+9}-3}{2x}$
        \item $\displaystyle\lim_{x\to 2}\frac{x^3-8}{x^2-4}$
    \end{enumerate}

    \item \textbf{[Subbab 4.2]} Dengan menggunakan definisi turunan, dapatkan $f'(x)$ dari $f(x)=\dfrac{1}{2x+1}$.
\end{enumerate}

\newpage

%%%%%%%%%%%%%%%%%%%%%%%%%%%%%%%%%%%%%%%%%%%%%%%%%%%%%%%%%%%%%%%%%%%%%%%%%%%%%%
% Kal-1 | Paket 2: subbab per nomor 1.2, 1.4, 2.2, 3.2, 4.2
% Kisi-kisi hari Selasa: subbab 1.2, 1.4, 2.1, 2.2, 3.2, 3.4, 4.2, 4.3
\kopsoal{Kal-1}{1-6}{Fisika, T. Kimia, T. Pangan}{Selasa, Kamis; 12.30-15.00}{8}
\begin{enumerate}
    \item \textbf{[Subbab 1.2]} Tentukan semua nilai $x$ yang memenuhi pertidaksamaan
    \begin{align*}
        |3x-1|\leq|x+3|.
    \end{align*}

    \item \textbf{[Subbab 1.4]} Diketahui titik $A(k,4)$, $B(1,1)$, dan $C(3,-2)$ terletak pada satu garis lurus.
    \begin{enumerate}
        \item Tentukan nilai $k$.
        \item Dapatkan persamaan garis tersebut.
    \end{enumerate}

    \item \textbf{[Subbab 2.2]} Diberikan fungsi $f(x)=\dfrac{1}{x+2}$ dan $g(x)=\sqrt{x-1}$.
    \begin{enumerate}
        \item Dapatkan domain $f$ dan $g$.
        \item Dapatkan $(g\circ f)(x)$ beserta domainnya.
    \end{enumerate}

    \item \textbf{[Subbab 3.2]} Hitunglah limit berikut.
    \begin{enumerate}
        \item $\displaystyle\lim_{t\to 2}\frac{t^2+t-6}{t^2-4}$
        \item $\displaystyle\lim_{x\to 5}\frac{\sqrt{2x-1}-3}{x-5}$
    \end{enumerate}

    \item \textbf{[Subbab 4.2]} Tunjukkan bahwa $f(x)=\begin{cases}
        3x-1, & x\leq 2\\
        x^2+1, & x>2
    \end{cases}$ adalah fungsi kontinu tetapi tidak dapat diturunkan di $x=2$.
\end{enumerate}

\newpage

%%%%%%%%%%%%%%%%%%%%%%%%%%%%%%%%%%%%%%%%%%%%%%%%%%%%%%%%%%%%%%%%%%%%%%%%%%%%%%
% Kal-2 | Paket 2: subbab per nomor 1.2, 1.4, 2.2, 3.2, 4.2
% Kisi-kisi hari Selasa: subbab 1.2, 1.4, 2.1, 2.2, 3.2, 3.4, 4.2, 4.3
\kopsoal{Kal-2}{7-12}{T. Lingkungan, T. Elektro, T. Biomedik}{Selasa, Kamis; 12.30-15.00}{8}
\begin{enumerate}
    \item \textbf{[Subbab 1.2]} Tentukan himpunan penyelesaian dari
    \begin{align*}
        \left|x^2-5\right|\leq 4.
    \end{align*}

    \item \textbf{[Subbab 1.4]} Diberikan garis $\ell_1: 3x-4y=12$ dan titik $P(1,4)$.
    \begin{enumerate}
        \item Dapatkan persamaan garis $\ell_2$ yang melalui $P$ dan tegak lurus dengan $\ell_1$.
        \item Tentukan titik potong garis $\ell_1$ dan $\ell_2$.
        \item Tentukan jarak titik $P$ ke garis $\ell_1$.
    \end{enumerate}

    \item \textbf{[Subbab 2.2]} Diberikan $f(x)=\sqrt{x+5}$ dan $(f\circ g)(x)=|x-1|$.
    \begin{enumerate}
        \item Dapatkan $g(x)$.
        \item Hitunglah $g(-1)$, kemudian periksa bahwa $(f\circ g)(-1)=|-1-1|$.
    \end{enumerate}

    \item \textbf{[Subbab 3.2]} Diketahui $\displaystyle\lim_{x\to 3}f(x)=4$ dan $\displaystyle\lim_{x\to 3}g(x)=-2$. Dapatkan
    \begin{enumerate}
        \item $\displaystyle\lim_{x\to 3}\big(3f(x)+g^2(x)\big)$
        \item $\displaystyle\lim_{x\to 3}\frac{\sqrt{f(x)+5}}{x\,g(x)}$
    \end{enumerate}

    \item \textbf{[Subbab 4.2]} Luas lingkaran berjari-jari $r$ adalah $L=\pi r^2$. Gunakan definisi turunan untuk mendapatkan $\dfrac{dL}{dr}$, kemudian hitung nilainya saat $r=3$ dan jelaskan artinya.
\end{enumerate}

\newpage

%%%%%%%%%%%%%%%%%%%%%%%%%%%%%%%%%%%%%%%%%%%%%%%%%%%%%%%%%%%%%%%%%%%%%%%%%%%%%%
% Kal-10 | Paket 2: subbab per nomor 1.2, 1.4, 2.1, 3.2, 4.3
% Kisi-kisi hari Selasa: subbab 1.2, 1.4, 2.1, 2.2, 3.2, 3.4, 4.2, 4.3
\kopsoal{Kal-10}{49-50}{Teknologi Kedokteran}{Selasa, Kamis; 12.30-15.00}{8}
\begin{enumerate}
    \item \textbf{[Subbab 1.2]} Tentukan himpunan penyelesaian dari
    \begin{align*}
        |x-1|<2|x+1|.
    \end{align*}

    \item \textbf{[Subbab 1.4]} Diberikan segitiga $ABC$ dengan $A(0,0)$, $B(6,2)$, dan $C(2,6)$.
    \begin{enumerate}
        \item Dapatkan persamaan garis tinggi dari titik $C$ (yaitu garis yang melalui $C$ dan tegak lurus $AB$).
        \item Tentukan jarak titik $C$ ke garis $AB$, kemudian hitung luas segitiga $ABC$.
    \end{enumerate}

    \item \textbf{[Subbab 2.1]} Diberikan $f(x)=\dfrac{\sqrt{x^2-2x-3}}{x-4}$.
    \begin{enumerate}
        \item Tentukan domain dari $f$.
        \item Hitunglah $f(5)$ dan $f(-2)$.
    \end{enumerate}

    \item \textbf{[Subbab 3.2]} Misalkan $f(x)=\begin{cases}
        \dfrac{x^2-4}{x-2}, & x<2\\
        x^2-2x+5, & x\geq 2
    \end{cases}$.
    \begin{enumerate}
        \item Dapatkan $\displaystyle\lim_{x\to 2^-}f(x)$ dan $\displaystyle\lim_{x\to 2^+}f(x)$.
        \item Apakah $\displaystyle\lim_{x\to 2}f(x)$ ada? Jelaskan.
    \end{enumerate}

    \item \textbf{[Subbab 4.3]} Tuliskan rumus turunan dari $\dfrac{f}{g}$, yaitu $\dfrac{d}{dx}\left(\dfrac{f(x)}{g(x)}\right)$.\\
    Selanjutnya, dapatkan turunan dari $\dfrac{f}{g}$ jika diketahui $f(x)=x+2$ dan $g(x)=x^2+1$.
\end{enumerate}

\newpage

%%%%%%%%%%%%%%%%%%%%%%%%%%%%%%%%%%%%%%%%%%%%%%%%%%%%%%%%%%%%%%%%%%%%%%%%%%%%%%
% Kal-iup-3 | Paket 2: subbab per nomor 1.2, 1.4, 2.2, 3.2, 4.2
% Kisi-kisi hari Selasa: subbab 1.2, 1.4, 2.1, 2.2, 3.2, 3.4, 4.2, 4.3
\kopsoaliup{Kal-iup-3}{63, 64, 70}{Civil, Geomatics, Geophysical Engineering, Naval Architecture (JD)}{Tuesday, Thursday; 15.30-17.10}{9}
\begin{enumerate}
    \item \textbf{[Subbab 1.2]} Find the solution set of the following inequality
    \begin{align*}
        |2x-5|\geq x+1.
    \end{align*}

    \item \textbf{[Subbab 1.4]} Find the equation of the line passing through the intersection point of the lines $x+y=5$ and $2x-y=1$, and perpendicular to the line $x+3y=6$.

    \item \textbf{[Subbab 2.2]} Given $f(x)=\dfrac{x}{x-1}$ and $g(x)=\dfrac{1}{x}$. Determine:
    \begin{enumerate}
        \item[a)] The domain of function $f$ and $g$;
        \item[b)] Function $(f\circ g)(x)$ and its domain.
    \end{enumerate}

    \item \textbf{[Subbab 3.2]} Calculate the following limits.
    \begin{enumerate}
        \item[a.] $\displaystyle\lim_{x\to 2}\frac{x^2-4}{x^2-3x+2}$
        \item[b.] $\displaystyle\lim_{x\to 9}\frac{x-9}{\sqrt{x}-3}$
    \end{enumerate}

    \item \textbf{[Subbab 4.2]} Using the definition of derivative, find $f'(x)$ of $f(x)=\dfrac{1}{x+2}$. Then, find the equation of the tangent line to the curve $y=f(x)$ at $x=0$.
\end{enumerate}

\newpage

\end{document}
